\documentclass[12pt, a4paper]{article}
\usepackage{../../../myconfig} % Gọi file cấu hình từ ngoài gốc vào

\begin{document}

% Tiêu đề đầu tài liệu
\titlebox{KIỂM TRA VẤN ĐÁP \& TỰ LUẬN}{TỔNG ÔN TRỌNG TÂM HÀM SỐ (CHUẨN THPTQG)}{MATH-STER: ÔN TẬP HÀM SỐ}

% =========================================================================
% DẠNG 1: ĐẠO HÀM NHANH
% =========================================================================
\dangbox{DẠNG 1}{Kỹ năng tính nhanh đạo hàm phân thức}

\questioncore{1}{Áp dụng công thức đạo hàm nhanh để tính đạo hàm của các hàm số sau:
\begin{enumerate}
    \item $y = \dfrac{2x+3}{x-4}$
    \item $y = \dfrac{2-x}{2x+1}$
    \item $y = \dfrac{x^2-2x+3}{x-1}$
    \item $y = \dfrac{2x^2+x-3}{x+2}$
\end{enumerate}
}

% =========================================================================
% DẠNG 2: ĐỒ THỊ f(x)
% =========================================================================
\newpage
\dangbox{DẠNG 2}{Khảo sát tính đơn điệu và cực trị từ đồ thị $y = f(x)$}

\questioncore{2}{Cho hàm số bậc ba $y = f(x)$ có đồ thị như hình vẽ:
\begin{center}
\begin{tikzpicture}[scale=0.9, >=stealth]
    \draw[->, line width=0.8pt] (-2.2,0) -- (2.5,0) node[below] {$x$};
    \draw[->, line width=0.8pt] (0,-2.8) -- (0,2.8) node[left] {$y$};
    \node[below left] at (0,0) {$O$};
    \draw[dashed] (1,0) -- (1,2) -- (0,2);
    \draw[dashed] (-1,0) -- (-1,-2) -- (0,-2);
    \node[below] at (1,0) {$1$};
    \node[above] at (-1,0) {$-1$};
    \node[left] at (0,2) {$2$};
    \node[right] at (0,-2) {$-2$};
    \draw[domain=-1.75:1.75, smooth, variable=\x, line width=1.2pt, deepblue] 
        plot ({\x}, {-(\x)^3 + 3*(\x)});
    \node[right, deepblue] at (1.3,1.8) {$y = f(x)$};
\end{tikzpicture}
\end{center}
Dựa vào đồ thị, hãy:
\begin{enumerate}
    \item Tìm các khoảng đồng biến và nghịch biến của hàm số $y = f(x)$.
    \item Xác định các điểm cực đại, điểm cực tiểu của hàm số và giá trị cực trị tương ứng.
\end{enumerate}
}

% =========================================================================
% DẠNG 3: ĐỒ THỊ f'(x)
% =========================================================================
\newpage
\dangbox{DẠNG 3}{Khảo sát tính đơn điệu và cực trị từ đồ thị $y = f'(x)$}

\questioncore{3}{Cho hàm số $y = f(x)$ có đạo hàm $f'(x)$ liên tục trên $\mathbb{R}$. Đồ thị của hàm số đạo hàm $y = f'(x)$ tiếp xúc với trục hoành tại điểm $x = -2$ và cắt trục hoành tại điểm $x = 1$ như hình vẽ bên dưới:
\begin{center}
\begin{tikzpicture}[scale=0.95, >=stealth]
    % Trục toạ độ
    \draw[->, line width=0.8pt] (-3.5,0) -- (2.5,0) node[below] {$x$};
    \draw[->, line width=0.8pt] (0,-3.2) -- (0,2.5) node[left] {$y$};
    \node[below left] at (0,0) {$O$};
    
    \node[above] at (-2,0) {$-2$};
    \node[below right] at (1,0) {$1$};
    
    % f'(x) = -0.3*(x+2)^2*(x-1) 
    % Tiếp xúc Ox tại x = -2 (nghiệm bội chẵn), cắt xuyên qua tại x = 1 (nghiệm đơn)
    \draw[domain=-3.1:1.6, smooth, variable=\x, line width=1.2pt, accentred] 
        plot ({\x}, {-0.3*((\x)+2)^2*((\x)-1)});
    \node[above right, accentred] at (0.2,1.5) {$y = f'(x)$};
    
    \filldraw[deepblue] (-2,0) circle (2pt);
    \filldraw[deepblue] (1,0) circle (2pt);
\end{tikzpicture}
\end{center}
Dựa vào đồ thị của $y = f'(x)$, hãy:
\begin{enumerate}
    \item Xác định các khoảng đồng biến và nghịch biến của hàm số $y = f(x)$.
    \item Tìm số điểm cực trị của hàm số $y = f(x)$ và chỉ rõ đâu là điểm cực đại, điểm cực tiểu (nếu có).
\end{enumerate}
}

% =========================================================================
% DẠNG 4: CÂU HỎI VẤN ĐÁP VỀ TIỆM CẬN
% =========================================================================
\newpage
\dangbox{DẠNG 4}{Câu hỏi kiểm tra vấn đáp về tiệm cận}

\questioncore{4}{Nêu định nghĩa và cách tìm đường tiệm cận đứng (TCĐ) của đồ thị hàm số.}

\questioncore{5}{Nêu định nghĩa và cách tìm đường tiệm cận ngang (TCN) của đồ thị hàm số.}

\questioncore{6}{Nêu định nghĩa và cách tìm đường tiệm cận xiên (TCX) của đồ thị hàm số.}

% =========================================================================
% DẠNG 5: BÀI TẬP TÌM TIỆM CẬN
% =========================================================================
\newpage
\dangbox{DẠNG 5}{Bài tập ví dụ tìm tiệm cận}

\questioncore{7}{Tìm tất cả các đường tiệm cận đứng và tiệm cận ngang của đồ thị hàm số:
$$y = \frac{x^2 - 3x + 2}{x^2 - 4}$$
}

\questioncore{8}{Tìm phương trình đường tiệm cận đứng và đường tiệm cận xiên của đồ thị hàm số:
$$y = \frac{x^2 - x + 1}{x - 1}$$
}

\questioncore{9}{Tìm các đường tiệm cận ngang và tiệm cận xiên của đồ thị hàm số:
$$y = \sqrt{x^2 + 2x} - x$$
}

% =========================================================================
% DẠNG 6: TÌM HÀM SỐ BẬC BA TỪ ĐỒ THỊ
% =========================================================================
\newpage
\dangbox{DẠNG 6}{Xác định hàm số bậc ba từ đồ thị}

\questioncore{10}{Đồ thị hàm số bậc ba $y = ax^3 + bx^2 + cx + d$ đi qua gốc tọa độ $O(0;0)$, đạt cực tiểu tại $A(-1;-2)$ và cực đại tại $B(1;2)$ như hình vẽ:
\begin{center}
\begin{tikzpicture}[scale=0.85, >=stealth]
    \draw[->, line width=0.8pt] (-2.5,0) -- (2.5,0) node[below] {$x$};
    \draw[->, line width=0.8pt] (0,-3.2) -- (0,3.2) node[left] {$y$};
    \node[below left] at (0,0) {$O$};
    \draw[dashed] (-1,0) -- (-1,-2) -- (0,-2);
    \draw[dashed] (1,0) -- (1,2) -- (0,2);
    \node[above] at (-1,0) {$-1$};
    \node[below] at (1,0) {$1$};
    \node[right] at (0,-2) {$-2$};
    \node[left] at (0,2) {$2$};
    \draw[domain=-1.7:1.7, smooth, variable=\x, line width=1.2pt, deepblue] 
        plot ({\x}, {-(\x)^3 + 3*(\x)});
    \filldraw[accentred] (-1,-2) circle (2pt) node[below] {$A(-1;-2)$};
    \filldraw[accentred] (1,2) circle (2pt) node[above] {$B(1;2)$};
\end{tikzpicture}
\end{center}
Hãy xác định phương trình của hàm số bậc ba đó.}

\questioncore{11}{Đồ thị hàm số bậc ba $y = f(x)$ tiếp xúc với trục hoành tại điểm $x = -1$, cắt trục hoành tại điểm $x = 2$ và cắt trục tung tại điểm $(0;2)$ như hình vẽ:
\begin{center}
\begin{tikzpicture}[scale=0.85, >=stealth]
    \draw[->, line width=0.8pt] (-2.2,0) -- (3,0) node[below] {$x$};
    \draw[->, line width=0.8pt] (0,-2.5) -- (0,3.5) node[left] {$y$};
    \node[below left] at (0,0) {$O$};
    \node[below left] at (-1,0) {$-1$};
    \node[below] at (2,0) {$2$};
    \node[right] at (0,2) {$2$};
    \draw[domain=-1.6:2.4, smooth, variable=\x, line width=1.2pt, deepblue] 
        plot ({\x}, {-((\x)+1)^2 * ((\x)-2)});
    \filldraw[accentred] (-1,0) circle (2pt);
    \filldraw[accentred] (2,0) circle (2pt);
    \filldraw[deepblue] (0,2) circle (2pt);
\end{tikzpicture}
\end{center}
Hãy xác định công thức của hàm số $y = f(x)$.}

% =========================================================================
% DẠNG 7: TÌM HÀM SỐ HỮU TỶ TỪ ĐỒ THỊ
% =========================================================================
\newpage
\dangbox{DẠNG 7}{Xác định hàm số hữu tỷ từ đồ thị}

\questioncore{12}{Cho hàm số bậc nhất trên bậc nhất $y = \dfrac{ax+b}{cx+d}$ có đồ thị như hình vẽ:
\begin{center}
\begin{tikzpicture}[scale=0.85, >=stealth]
    \draw[->, line width=0.8pt] (-2.5,0) -- (4,0) node[below] {$x$};
    \draw[->, line width=0.8pt] (0,-2.5) -- (0,4) node[left] {$y$};
    \node[below left] at (0,0) {$O$};
    
    \draw[dashed, line width=0.9pt, accentred] (1,-2.5) -- (1,4) node[above] {$x=1$};
    \draw[dashed, line width=0.9pt, accentred] (-2.5,1) -- (4,1) node[right] {$y=1$};
    
    \draw[domain=-2.5:0.75, smooth, variable=\x, line width=1.2pt, deepblue] 
        plot ({\x}, {((\x)-2)/((\x)-1)});
    \draw[domain=1.25:3.8, smooth, variable=\x, line width=1.2pt, deepblue] 
        plot ({\x}, {((\x)-2)/((\x)-1)});
        
    \filldraw[deepblue] (0,2) circle (2pt) node[left] {$2$};
    \filldraw[deepblue] (2,0) circle (2pt) node[below] {$2$};
\end{tikzpicture}
\end{center}
\begin{enumerate}
    \item Xác định phương trình các đường tiệm cận đứng, tiệm cận ngang và xét dấu của biểu thức $ad - bc$.
    \item Tìm các hệ số $a, b, c$ của hàm số dạng $y = \dfrac{ax+2}{bx+c}$ biết đồ thị đi qua các điểm và có tiệm cận như hình.
\end{enumerate}
}

\questioncore{13}{Cho hàm số bậc hai trên bậc nhất $y = \dfrac{x^2 - x + 1}{x - 1}$ có đồ thị như hình vẽ:
\begin{center}
\begin{tikzpicture}[scale=0.85, >=stealth]
    \draw[->, line width=0.8pt] (-2.5,0) -- (4.5,0) node[below] {$x$};
    \draw[->, line width=0.8pt] (0,-3) -- (0,4.5) node[left] {$y$};
    \node[below left] at (0,0) {$O$};
    
    \draw[dashed, line width=0.9pt, accentred] (1,-3) -- (1,4.5) node[above] {$x=1$};
    \draw[dashed, line width=0.9pt, amberorange] (-2.5,-2.5) -- (4,4) node[right] {$y=x$};
    
    \draw[domain=-2.2:0.72, smooth, variable=\x, line width=1.2pt, deepblue] 
        plot ({\x}, {((\x)^2 - (\x) + 1)/((\x) - 1)});
    \draw[domain=1.28:3.8, smooth, variable=\x, line width=1.2pt, deepblue] 
        plot ({\x}, {((\x)^2 - (\x) + 1)/((\x) - 1)});
        
    \filldraw[accentred] (0,-1) circle (2pt) node[left] {$(0;-1)$};
    \filldraw[accentred] (2,3) circle (2pt) node[right] {$(2;3)$};
\end{tikzpicture}
\end{center}
Hãy xác định phương trình tiệm cận đứng, tiệm cận xiên và tọa độ hai điểm cực trị của đồ thị hàm số.}

% =========================================================================
% DẠNG 8: BÀI TOÁN THỰC TẾ & TỐI ƯU HÓA (ĐỀ MỚI HOÀN TOÀN)
% =========================================================================
\newpage
\dangbox{DẠNG 8}{Ứng dụng đạo hàm giải quyết bài toán thực tế}

\questioncore{14}{Sau khi bệnh nhân được tiêm một liều kháng sinh, nồng độ thuốc $C(t)$ (đơn vị: mg/lít) trong máu sau $t$ giờ được xác định bởi hàm số:
$$C(t) = -\frac{1}{18}t^3 + \frac{1}{2}t^2 + 8t \quad (0 \le t \le 10)$$
Hỏi tốc độ tăng nồng độ thuốc trong máu của bệnh nhân đạt giá trị lớn nhất sau khi tiêm bao nhiêu giờ?}

\vspace{3cm}

\questioncore{15}{Một hợp tác xã nông nghiệp dự định dùng $240\text{ m}$ hàng rào lưới thép để ngăn một khu đất trống tựa lưng vào bức tường rào bê tông có sẵn thành 3 ô chuồng hình chữ nhật giống nhau (như hình vẽ, bờ tường bê tông không cần rào):
\begin{center}
\begin{tikzpicture}[scale=0.85, >=stealth]
    \draw[line width=2pt, gray!80] (-0.8,0) -- (6.8,0) node[right, black] {Bức tường bê tông};
    \draw[line width=1.2pt, deepblue] (0,0) -- (0,2.2) -- (6,2.2) -- (6,0);
    \draw[line width=1.2pt, deepblue] (2,0) -- (2,2.2);
    \draw[line width=1.2pt, deepblue] (4,0) -- (4,2.2);
    \node at (1,1.1) {Ô 1};
    \node at (3,1.1) {Ô 2};
    \node at (5,1.1) {Ô 3};
    \draw[<->] (-0.3,0) -- (-0.3,2.2) node[midway, left] {$x$};
    \draw[<->] (0,2.5) -- (6,2.5) node[midway, above] {$y$};
\end{tikzpicture}
\end{center}
Hỏi tổng diện tích lớn nhất của 3 ô chuồng thu được là bao nhiêu mét vuông?}
\vspace{3cm}

\questioncore{16}{Một xưởng cơ khí chính xác sản xuất linh kiện vi mạch. Nếu mỗi tháng xuất xưởng $x$ trăm linh kiện thì doanh thu nhận được ước tính theo hàm $R(x) = -0{,}05x^2 + 650x$ (triệu đồng). Chi phí sản xuất trung bình cho mỗi trăm linh kiện được cho bởi công thức:
$$C_{tb}(x) = \frac{32000}{x} + 250 \quad (\text{triệu đồng})$$
Hỏi xưởng cơ khí cần sản xuất bao nhiêu trăm linh kiện trong một tháng để đạt được lợi nhuận tối đa?}

\vspace{3cm}

\newpage
\questioncore{17}{Để cấp điện cho một cụm dân cư trên đảo $D$, người ta kéo đường cáp tải điện từ nhà máy phát điện $M$ đặt trên bờ biển. Khoảng cách vuông góc ngắn nhất từ đảo $D$ vào đất liền tại điểm $H$ là $DH = 9\text{ km}$, khoảng cách dọc theo bờ biển từ $H$ đến nhà máy $M$ là $HM = 27\text{ km}$. Chi phí kéo dây cáp trên đất liền từ nhà máy $M$ đến trạm trung chuyển $K$ ($K$ nằm giữa $H$ và $M$) là $400$ triệu đồng/km, chi phí ngầm hóa đáy biển từ $K$ ra đảo $D$ là $500$ triệu đồng/km:
\begin{center}
\begin{tikzpicture}[scale=0.22, >=stealth]
    \coordinate (H) at (0,0);
    \coordinate (D) at (0,9);
    \coordinate (K) at (12,0);
    \coordinate (M) at (27,0);
    
    \draw[line width=1.2pt] (H) -- (M) node[below] {$M$};
    \draw[line width=1.2pt] (H) -- (D) node[above] {$D$};
    \draw[line width=1.2pt, dashed, deepblue] (D) -- (K);
    
    \node[below] at (H) {$H$};
    \node[below] at (K) {$K$};
    \node[left] at (0,4.5) {$9\text{ km}$};
    \node[above] at (6,4.5) {Dưới đáy biển};
    \node[below] at (19.5,0) {Trên đất liền};
    \draw[<->] (0,-3.5) -- (27,-3.5) node[midway, below] {$27\text{ km}$};
\end{tikzpicture}
\end{center}
Tìm chi phí thấp nhất để hoàn thành công trình đường dây tải điện trên (tính bằng tỉ đồng).}

\vspace{3cm}

\questioncore{18}{Một phần đường chạy thử nghiệm của phương tiện lượn trên đệm từ được mô phỏng bởi đồ thị hàm số bậc ba $y = f(x)$ như hình vẽ:
\begin{center}
\begin{tikzpicture}[scale=0.09, >=stealth]
    % Trục tọa độ
    \draw[->, line width=0.8pt] (-5,0) -- (105,0) node[below] {$x$};
    \draw[->, line width=0.8pt] (0,-4) -- (0,45) node[left] {$y$};
    \node[below left] at (0,0) {$O$};
    
    % Đường cao độ chuẩn y = 20
    \draw[dashed, line width=0.6pt, black!40] (0,20) -- (90,20);
    \node[left] at (0,20) {$20$};
    
    % Vẽ nhánh đồ thị: f(x) = -(1/3500)*x*(x-60)*(x-90) + 20
    \draw[domain=0:96, smooth, variable=\x, line width=1.2pt, deepblue] 
        plot ({\x}, {-(1/3500)*(\x)*((\x)-60)*((\x)-90) + 20});
        
    % Các điểm giao trên đường y = 20
    \filldraw[accentred] (0,20) circle (16pt) node[above left] {$A(0;20)$};
    \filldraw[accentred] (60,20) circle (16pt) node[above left] {$C(60;20)$};
    \filldraw[accentred] (90,20) circle (16pt) node[above right] {$D(90;20)$};
    
    % Điểm trũng nhất tại x = 20, y = 4
    \draw[dashed, line width=0.6pt, black!50] (20,0) -- (20,4) -- (0,4);
    \node[below] at (20,0) {$20$};
    \node[left] at (0,4) {$4$};
    \filldraw[deepblue] (20,4) circle (16pt) node[below right] {\small Điểm võng};
\end{tikzpicture}
\end{center}
Biết đường chạy xuất phát từ điểm $A(0;20)$, đi qua các điểm $C(60;20)$, $D(90;20)$ và đạt độ cao thấp nhất là $4\text{ m}$ tại vị trí có hoành độ $x = 20$. Hãy xác định độ cao lớn nhất mà phương tiện đạt được trên quỹ đạo chạy.}
\vspace{3cm}
\newpage

\questioncore{19}{Một bể chứa nước chữa cháy có dạng hình lăng trụ đứng với đáy là tam giác vuông cân tại đỉnh đáy, chiều dài máng $L = 8\text{ m}$ và độ sâu tối đa $H = 2\text{ m}$ (mép trên miệng máng rộng $W = 4\text{ m}$) như hình vẽ:
\begin{center}
\begin{tikzpicture}[scale=0.9, >=stealth]
    % Đỉnh đáy và các điểm mặt trước
    \coordinate (O) at (0,0);
    \coordinate (A) at (-2,2);
    \coordinate (B) at (2,2);
    
    % Chiều sâu không gian lăng trụ (vector trượt)
    \coordinate (shift) at (4,1.5);
    \coordinate (O2) at ($(O)+(shift)$);
    \coordinate (A2) at ($(A)+(shift)$);
    \coordinate (B2) at ($(B)+(shift)$);
    
    % Mặt nước ở độ cao h
    \coordinate (M) at (-1.2,1.2);
    \coordinate (N) at (1.2,1.2);
    \coordinate (M2) at ($(M)+(shift)$);
    \coordinate (N2) at ($(N)+(shift)$);
    
    % Vùng nước
    \fill[cyan!20, opacity=0.8] (O) -- (M) -- (M2) -- (O2) -- cycle;
    \fill[cyan!30, opacity=0.8] (O) -- (N) -- (N2) -- (O2) -- cycle;
    \fill[cyan!40, opacity=0.8] (M) -- (N) -- (N2) -- (M2) -- cycle;
    
    % Viền bể lăng trụ
    \draw[line width=1.1pt, deepblue] (A) -- (O) -- (B) -- cycle;
    \draw[line width=1.1pt, deepblue] (A) -- (A2) -- (B2) -- (B);
    \draw[line width=1.1pt, deepblue] (O) -- (O2) -- (B2);
    \draw[dashed, line width=0.8pt, deepblue!60] (O2) -- (A2);
    
    % Kích thước chiều dài L = 8m
    \draw[<->, line width=0.7pt] ($(B)+(0.2,-0.1)$) -- ($(B2)+(0.2,-0.1)$) node[midway, below right] {$L = 8\text{ m}$};
    
    % Mặt trước: chiều cao h và độ sâu H
    \draw[dashed, line width=0.6pt] (0,2) -- (O);
    \draw[dashed, line width=0.6pt] (M) -- (N);
    \draw[<->, line width=0.6pt] (-2.3,0) -- (-2.3,2) node[midway, left] {$H = 2\text{ m}$};
    \draw[<->, line width=0.6pt] (0.3,0) -- (0.3,1.2) node[midway, right] {$h$};
    \draw[<->, line width=0.6pt] (-2,2.2) -- (2,2.2) node[midway, above] {$W = 4\text{ m}$};
\end{tikzpicture}
\end{center}
Ban đầu bể không có nước, người ta bơm nước vào bể với tốc độ không đổi $q = 1{,}6\text{ m}^3/\text{phút}$.
\begin{enumerate}
    \item Sau 5 phút bơm, mực nước $h$ trong bể cao bao nhiêu mét?
    \item Tại thời điểm mực nước trong bể đạt độ cao $h = 1\text{ m}$, tốc độ dâng lên của mực nước là bao nhiêu mét/phút?
\end{enumerate}
}

\end{document}
```<FollowUp>
Bạn có muốn tôi chèn thêm các dòng kẻ ô ly hoặc dòng chấm lửng (theo macro có sẵn trong myconfig) dưới mỗi câu để học sinh có không gian viết lời giải trực tiếp vào đề không?
</FollowUp>